\documentclass{article}
\usepackage[UTF8]{ctex}
\usepackage{amsmath,amssymb,amsfonts}
\usepackage{enumitem}
\usepackage{geometry}
\geometry{a4paper, margin=1in}
% === 错题本排版优化（不参与题目标识） ===
\setlength{\parindent}{0pt}
\setlength{\parskip}{0.35em plus 0.10em minus 0.08em}
\setlength{\emergencystretch}{3em}
\setlist{itemsep=0.25em, topsep=0.40em, parsep=0pt}
\allowdisplaybreaks[2]
% === 错题本排版优化结束 ===
\begin{document}
\section*{西南大学2024年高等代数真题}
\section*{56．西南大学2024年高等代数试卷}

\subsection{题目1}
1．（20 分）计算 \(n\) 阶行列式

\[
D_{n}=\left|\begin{array}{cccc}
a+x_{1} & a+x_{1}^{2} & \cdots & a+x_{1}^{n} \\
a+x_{2} & a+x_{2}^{2} & \cdots & a+x_{2}^{n} \\
\vdots & \vdots & & \vdots \\
a+x_{n} & a+x_{n}^{2} & \cdots & a+x_{n}^{n}
\end{array}\right| .
\]

\subsection{题目2}
2．（20 分）设矩阵 \(A=\left(\begin{array}{lll}1 & 4 & 1 \\ 1 & 1 & 2 \\ 0 & 0 & 5\end{array}\right)\) ，求 \(A^{2024}\) ．

\subsection{题目3}
3．（20 分）设 \(A\) 是一个二阶实对称矩阵，\(A\) 的特征值为 \(\lambda_{1}=2, \lambda_{2}=3\) ．并设 \(\alpha=\binom{1}{1}\) 是属于特征值 \(\lambda_{1}=2\) 的特征向量，求矩阵 \(A\) 。

\subsection{题目4}
4．（20 分）请问 \(a, b\) 取什么值时，线性方程组

\[
\left\{\begin{array}{l}
x_{1}+x_{2}+x_{3}+x_{4}+x_{5}=1 \\
2 x_{1}+x_{2}-4 x_{5}=a \\
x_{1}+2 x_{2}+3 x_{3}+3 x_{4}+7 x_{5}=4 \\
5 x_{1}+4 x_{2}+3 x_{3}+3 x_{4}-x_{5}=b
\end{array}\right.
\]

有解，并求其解。

\subsection{题目5}
5．（20 分）设

\[
f(x)=x^{3}+a x^{2}+b x+c
\]

是一个整系数多项式，且 \(a c+b c\) 是奇数．证明：\(f(x)\) 在有理数域上不可约．

\subsection{题目6}
6．（20 分）设 \(f\left(x_{1}, x_{2}, \cdots, x_{n}\right)\) 是一个 \(n\) 元实二次型．证明：存在 \(n\) 元正定二次型 \(g\left(x_{1}, x_{2}, \cdots, x_{n}\right)\) 和\(n\) 元负定二次型 \(h\left(x_{1}, x_{2}, \cdots, x_{n}\right)\) 使得

\[
f\left(x_{1}, x_{2}, \cdots, x_{n}\right)=g\left(x_{1}, x_{2}, \cdots, x_{n}\right)+h\left(x_{1}, x_{2}, \cdots, x_{n}\right) .
\]

\subsection{题目7}
7．（15 分）设 \(\sigma\) 为 \(n\) 维线性空间 \(V\) 上的线性变换，\(\lambda\) 为 \(\sigma\) 的一个特征根，\(k\) 为 \(\lambda\) 的代数重数．记

\[
W=\operatorname{Ker}\left(\lambda I_{V}-\sigma\right)^{k} .
\]

其中 \(I_{V}\) 是 \(V\) 上的恒等变换。证明：\(W\) 的维数恰为 \(k\) 。

\subsection{题目8}
8．（15 分）解答如下问题：\\
（1）设 \(V\) 为数域 \(F\) 上的线性空间，\(W\) 为集合，\(f\) 为 \(V\) 到 \(W\) 的双射．证明：\(W\) 也可以做成数域 \(F\) 上的线性空间。\\
（2）自然数集是否可以构成有理数域上的 2 维线性空间？请说明理由。
\end{document}
